\title{
Lärarenkät för det nationella provet i engelska årskurs 6, 2023/2024
}
\section*{GÖTEBORGS}
UNIVERSITET
\section*{Lärarenkät Eng åk6 VT 2024}
\section*{Sammanställningen av denna webbenkät grundar sig på alla inkomna enkatsvar, 419 stycken. Här redovisas alla frågor med svarsalternativ. Siffrorna är avrundade till heltal.}
Enkäten ger också flera möjligheter till att lämna öppna kommentarer och synpunkter. Dessa tas upp i resultatrapporten för det nationella provet i engelska årskurs 6, läsåret 2023/2024.
Eftersom antalet besvarade enkäter är relativt lägt och självselekterat bör resultaten tolkas med försiktighet.
\section*{Skolform}
\(100 \%\) Grundskola
\(0 \%\) Specialskola
\(0 \%\) Om annat, specificera
\section*{Provets genomförande}
I vilken utsträckning stämmer följande påståenden in på din uppfattning?
\begin{tabular}{|c|c|c|c|c|}
\hline & \begin{tabular}{l} 
Instämmer \\
helt
\end{tabular} & \begin{tabular}{l} 
Instämmer till \\
stor del
\end{tabular} & \begin{tabular}{l} 
Instämme till \\
viss del
\end{tabular} & \begin{tabular}{l} 
Instämmer \\
inte alls
\end{tabular} \\
\hline Provet som helhet är bra. & \(61 \%\) & \(35 \%\) & \(4 \%\) & \(0 \%\) \\
\hline \begin{tabular}{l} 
Provet som helhet är ett stöd för \\
betygssättningen.
\end{tabular} & \(64 \%\) & \(27 \%\) & \(7 \%\) & \(2 \%\) \\
\hline \begin{tabular}{l} 
Bedömningsanvisningarna har varit ett \\
bra stöd i bedömningen av elevernas \\
prestationer i provet.
\end{tabular} & \(53 \%\) & \(36 \%\) & \(9 \%\) & \(2 \%\) \\
\hline \begin{tabular}{l}  
Provet som helhet bidrar till att \\
konkretisera kursplanen
\end{tabular} & \(54 \%\) & \(37 \%\) & \(9 \%\) & \(1 \%\) \\
\hline
\end{tabular}
På vilket sätt har du förberett eleverna inför genomförandet av provet? Flera alternativ möjliga.
Genom att:
\(71 \% \quad\) informera om brev till elever respektive brev till vårdnadshavare
\(91 \% \quad\) använda information från hätet Lärarinformation
\(95 \% \quad\) låta eleverna genomföra uppgifter ur tidigare nationella prov
\(61 \% \quad\) använda Skolverkets bedömningsstöd
\(12 \% \quad\) annat, nämligen
Anser du att provresultaten ligger i linje med elevernas övriga prestationer under läsåret?
\(87 \% \quad\) Ja, för flertalet
\(10 \% \quad\) Ja, för cirka hälften
\(2 \% \quad\) Ja, men bara för ett fåtal
\(1 \% \quad\) Nej

Delprov A, Focus: Speaking
Anser du att Delprov A (Focus: Speaking) ger eleverna möjlighet att visa sin förmåga att formulera sig och kommunicera på engelska i tal?
\(93 \% \quad\) Ja
\(7 \% \quad\) Till viss del
\(0 \% \quad\) Nej
I relation till kursplanen, var uppgiften ...
\(70 \% \quad\) mycket lämplig
\(30 \% \quad\) lämplig
\(0 \% \quad\) mindre lämplig
I relation till kursplanen, var uppgiften ...
\(1 \% \quad\) svår
\(92 \% \quad\) lagom svår/lagom lätt
\(6 \% \quad\) lätt
Uppgiften fungerade ...
\(84 \%\) bra
\(15 \%\) ganska bra
\(0 \% \quad\) mindre bra
Vid bedömningen gav de kommenterade exemplen på elevprestationer ...
\(34 \% \quad\) mycket gott stöd
\(57 \% \quad\) tämligen gott stöd
\(9 \% \quad\) visst stöd
\(0 \% \quad\) ringa stöd
Hur sker bedömningen av delprov A?
\(34 \% \quad\) Samtliga elevprestationer sambedöms
\(11 \% \quad\) Många elevprestationer sambedöms
\(24 \% \quad\) Vissa elevprestationer sambedöms
\(29 \% \quad\) Samtliga elevprestationer bedöms av mig ensam
\(1 \% \quad\) Samtliga elevprestationer bedöms av annan lärare
\(1 \% \quad\) Annat, nämligen
Bedömningen av de inspelade elevprestationerna är för respektive betyg ...
\begin{tabular}{llll} 
& för sträng & rimlig & för mild \\
E & \(1 \%\) & \(82 \%\) & \(17 \%\) \\
D & \(0 \%\) & \(93 \%\) & \(7 \%\) \\
C & \(0 \%\) & \(93 \%\) & \(7 \%\) \\
B & \(0 \%\) & \(89 \%\) & \(11 \%\) \\
A & \(0 \%\) & \(85 \%\) & \(15 \%\)
\end{tabular}
Uppgiften genomfördes ... (Flera alternativ möjliga.)
\(65 \% \quad\) med inspelning
\(36 \% \quad\) utan inspelning

\title{
Delprov B 1, Focus: Reading
}
Anser du att Delprov B 1 (Focus: Reading) ger ett tillförlitligt mått på dina elevers förmåga att förstå och tolka innehållet i olika slags texter på engelska?
\(90 \% \quad \mathrm{Ja}\)
\(10 \% \quad\) Till viss del
\(1 \% \quad\) Nej
I relation till kursplanen, var respektive uppgift ...
mycket lämplig
\(56 \%\)
\(55 \%\)
\(55 \%\)
\(56 \%\)
lämplig
\(44 \%\)
\(44 \%\)
\(45 \%\)
\(43 \%\)
mindre lämplig
\(0 \%\)
\(1 \%\)
\(0 \%\)
\(1 \%\)
I relation till kursplanen, var respektive uppgift ...
\begin{tabular}{l|l|l|l} 
& svår & lagom svår/lagom lättt & lätt \\
\hline Uppgift B1a & \(2 \%\) & \(81 \%\) & \(17 \%\) \\
\hline Uppift B1b & \(4 \%\) & \(87 \%\) & \(9 \%\) \\
\hline Uppift B1c & \(3 \%\) & \(87 \%\) & \(10 \%\) \\
\hline Uppift B1d & \(6 \%\) & \(87 \%\) & \(7 \%\)
\end{tabular}
\section*{Delprov B 2, Focus: Listening}
Anser du att Delprov B 2 (Focus: Listening) ger ett tillförlitligt mått på dina elevers förmågra att förstå och tolka innehållet i talad engelska?
\(88 \% \quad\) Ja
\(11 \% \quad\) Till viss del
\(1 \% \quad\) Nej
I ration till ämnesplanen, var respektive uppgift ...
\begin{tabular}{llll} 
mycket lämplig & lämplig & mindre lämplig \\
Uppgift B2a & \(50 \%\) & \(48 \%\) & \(2 \%\) \\
Uppgift B2b & \(49 \%\) & \(49 \%\) & \(2 \%\)
\end{tabular}
I relation till ämnesplanen, var respektive uppgift ...
\begin{tabline & svår & lagom svår/lagom lättt & lätt \\
Uppgift B2a & \(4 \%\) & \(84 \%\) & \(12 \%\) \\
Uppift B2b & \(4 \%\) & \(87 \%\) & \(8 \%\)
\end{tabular}

\title{
Delprov B
}
Anser du att poänggränserna är rimliga?
(Bedömningsanvisningar 2, kapitel 4)
\(88 \% \quad\) Ja
\(9 \% \quad\) Till viss del
\(3 \% \quad\) Nej
Bedömningsanvisningarna till Delprov B ...
\(98 \%\) fungerade bra
\(2 \% \quad\) fungerade mindre bra
Hur sker bedömningen av delprov B?
\(26 \% \quad\) Samtliga elevprestationer sambedöms
\(14 \% \quad\) Många elevprestationer sambedöms
\(29 \% \quad\) Vissa elevprestationer sambedöms
\(26 \% \quad\) Samtliga elevprestationer bedöms av mig ensam
\(3 \% \quad\) Samtliga elevprestationer bedöms av annan lärare
\(2 \% \quad\) Annat, nämligen
\section*{Delprov C, Focus: Writing}
Anser du att Delprov C (Focus: Writing) ger eleverna möjlighet att visa sin förmåga att formulera sig och kommunicera på engelska i skrift?
\(76 \% \quad\) Ja
\(23 \% \quad\) Till viss del
\(1 \% \quad\) Nej
I relation till ämnesplanen, var uppgiften ...
\(37 \% \quad\) mycket lämplig
\(53 \% \quad\) lämplig
\(10 \% \quad\) mindre lämplig
I relation till kursplanen, var uppgiften ...
\(23 \% \quad\) svår
\(73 \% \quad\) lagom svår/lagom lätt
\(4 \% \quad\) lätt
Vid bedömningen gav de kommenterade exemplen på elevprestationer ...
\(31 \% \quad\) mycket gott stöd
\(51 \% \quad\) tämligen gott stöd
\(16 \% \quad\) visst stöd
\(2 \% \quad\) ringa stöd

Bedömningen av elevprestationerna är för respektive betyg ...
\begin{tabular}{llll} 
& för sträng & rimlig & för mild \\
E & \(1 \%\) & \(71 \%\) & \(28 \%\) \\
D & \(0 \%\) & \(80 \%\) & \(20 \%\) \\
C & \(0 \%\) & \(75 \%\) & \(25 \%\) \\
B & \(1 \%\) & \(73 \%\) & \(26 \%\) \\
A & \(1 \%\) & \(69 \%\) & \(30 \%\)
\end{tabular}
Hur sker bedömningen av delprov C?
\begin{tabular}{ll} 
63\% & Samtliga elevprestationer sambedöms \\
13\% & Många elevprestationer sambedöms \\
10\% & Vissa elevprestationer sambedöms \\
9\% & Samtliga elevprestationer bedöms av mig ensam \\
2\% & Samtliga elevprestationer bedöms av annan lärare \\
\(2 \%\) & Annat, nämligen
\end{tabular}
\section*{Provet som helhet}
I vilken utsträckning anser du att ämnesprovet som helhet ger resultat som stämmer överens med dina egna bedömningar av enskilda elevers kunskaper?
\(50 \% \quad\) Mycket stor
\(46 \%\) Ganska stor
\(4 \%\) Ganska liten
\(0 \% \quad\) Mycket liten
Har du svarat på den här enkäten själv, eller tillsammans med kollegor?
\(89 \%\) Jag har svarat på frågorna själv
\(10 \%\) Flera kollegor har lämnat ett gemensamt svar på frågorna
\(1 \% \quad\) På annat sätt, nämligen